\chapter{Resultados e discussão}\label{chp:resultados}

Este capítulo reserva a estrutura da análise final. Como os artefatos
consolidados das Fases 8 a 10 ainda serão atualizados, os resultados numéricos
permanecem como campos explícitos de preenchimento. Essa escolha evita atribuir
desempenho, robustez ou superioridade clínica antes da execução final
auditável. O objetivo desta versão é garantir que, quando os resultados forem
gerados, eles tenham lugar claro no texto, nas tabelas e nas figuras.

\section{Estatísticas das coortes}\label{sec:res_coortes}

A primeira análise quantitativa deverá descrever o volume de dados processado em
cada base: número de lâminas, identificadores de paciente, \textit{patches}
aceitos, \textit{patches} rejeitados e distribuição entre treino, validação e
teste. A \Tabela{tab:coortes_placeholder} apresenta a estrutura mínima esperada.

\begin{table}[!htb]
\centering
\caption{Estrutura provisória para estatísticas de coorte e particionamento.}
\label{tab:coortes_placeholder}
\begin{tabular}{p{3.2cm}p{1.7cm}p{1.6cm}p{2.2cm}p{3.6cm}}
\toprule
Base & WSIs & IDs & \textit{Patches} & Observação \\
\midrule
CAMELYON16 & -- & -- & -- & A preencher com manifesto final. \\
CATCH & -- & -- & -- & A preencher com manifesto final. \\
DiagSet & -- & -- & -- & A preencher com manifesto final. \\
HiESD & -- & -- & -- & A preencher com manifesto final. \\
\bottomrule
\end{tabular}
\end{table}

Além das contagens globais, a versão final deve relatar quantas amostras foram
removidas por baixo conteúdo tecidual, por ambiguidade de anotação, por falha de
leitura ou por regra de limpeza por grafos. Esses números são necessários para
avaliar se a curadoria alterou substancialmente a distribuição dos dados.

\section{Extração, limpeza e particionamento}\label{sec:res_preprocessamento}

Os resultados de pré-processamento deverão mostrar se o fluxo preservou os
contratos definidos na metodologia: alinhamento entre imagem e máscara,
separação de pacientes, existência de artefatos de normalização e ausência de
linhas inválidas. A discussão deve separar três dimensões. A primeira é
operacional, relacionada a tempo de execução, tamanho dos HDF5 e possibilidade de
retomada. A segunda é científica, relacionada a vazamento, rótulos ambíguos e
qualidade da supervisão. A terceira é prática, relacionada ao custo de disco e ao
uso de armazenamento compacto de treino.

\begin{figure}[!htb]
    \centering
    \fbox{\begin{minipage}{0.82\textwidth}
        \centering
        \vspace{2.5em}
        Figura qualitativa pendente: exemplos de \textit{patches} aceitos,
        rejeitados por artefato e rejeitados por ambiguidade.
        \vspace{2.5em}
    \end{minipage}}
    \caption[Placeholder para exemplos de curadoria]{Local reservado para exemplos qualitativos da curadoria de dados. As imagens finais devem ser extraídas dos artefatos atuais, não do rascunho antigo.}
    \label{fig:placeholder-curadoria}
\end{figure}

\section{Triagem e treinamento de modelos}\label{sec:res_treinamento}

A triagem da Fase 7 deverá ser relatada como seleção de configurações, não como
desempenho final. A Fase 8 deverá apresentar os modelos efetivamente treinados,
as configurações de perda e normalização, o melhor checkpoint selecionado por
validação e as métricas de validação correspondentes. A \Tabela{tab:modelos_placeholder}
define uma estrutura inicial para esse relato.

\begin{table}[!htb]
\centering
\caption{Estrutura provisória para modelos treinados na Fase 8.}
\label{tab:modelos_placeholder}
\begin{tabular}{p{2.5cm}p{2.8cm}p{2.6cm}p{2.4cm}p{2.4cm}}
\toprule
Arquitetura & Encoder & Normalização & Critério de seleção & Status \\
\midrule
A preencher & A preencher & A preencher & AUPRC validação & Pendente \\
A preencher & A preencher & A preencher & AUPRC validação & Pendente \\
A preencher & A preencher & A preencher & AUPRC validação & Pendente \\
\bottomrule
\end{tabular}
\end{table}

As curvas de treinamento deverão ser discutidas com cautela. Uma curva de perda
decrescente indica otimização, mas não garante generalização. Divergência,
colapso de predição, métricas não finitas e instabilidade entre épocas devem ser
reportadas como evidência metodológica, mesmo quando não produzirem o melhor
modelo.

\begin{figure}[!htb]
    \centering
    \fbox{\begin{minipage}{0.82\textwidth}
        \centering
        \vspace{2.5em}
        Figura pendente: curvas finais de perda, AUPRC, AUROC e MCC por época.
        \vspace{2.5em}
    \end{minipage}}
    \caption[Placeholder para curvas de treinamento]{Local reservado para curvas finais de treinamento e validação exportadas dos logs atuais.}
    \label{fig:placeholder-curvas}
\end{figure}

\section{Otimização do ensemble e inferência final}\label{sec:res_ensemble}

Os resultados da Fase 9 deverão descrever a composição do \textit{ensemble}, os
pesos atribuídos a cada modelo, o limiar calibrado e as assinaturas de
compatibilidade verificadas antes da inferência. A Fase 10 deverá apresentar
somente a avaliação da receita congelada no teste.

\begin{table}[!htb]
\centering
\caption{Estrutura provisória para métricas finais de inferência.}
\label{tab:metricas_finais_placeholder}
\small
\begin{tabular}{p{3.1cm}p{1.5cm}p{1.5cm}p{2.7cm}p{2.8cm}}
\toprule
Agregação & Dice & IoU & Sensibilidade & Especificidade \\
\midrule
Micro & -- & -- & -- & -- \\
Macro por paciente & -- & -- & -- & -- \\
Bootstrap por paciente & -- & -- & -- & -- \\
\bottomrule
\end{tabular}
\end{table}

\begin{figure}[!htb]
    \centering
    \fbox{\begin{minipage}{0.82\textwidth}
        \centering
        \vspace{2.5em}
        Figura qualitativa pendente: WSI ou \textit{patches} com máscara real,
        probabilidade do ensemble e máscara binarizada.
        \vspace{2.5em}
    \end{minipage}}
    \caption[Placeholder para inferência qualitativa]{Local reservado para exemplos qualitativos da inferência final.}
    \label{fig:placeholder-inferencia}
\end{figure}

\section{Discussão provisória}\label{sec:discussao_provisoria}

A discussão final deverá responder a quatro perguntas. Primeiro, o fluxo
preservou isolamento por paciente e rastreabilidade suficiente para sustentar as
métricas? Segundo, as etapas de controle de qualidade removeram ruído sem
eliminar de forma desproporcional regiões relevantes? Terceiro, a normalização
de cor e a amostragem inteligente alteraram estabilidade de treinamento ou
desempenho? Quarto, o \textit{ensemble} congelado melhorou a avaliação de teste
em relação aos modelos individuais?

Enquanto essas perguntas não forem respondidas pelos artefatos finais, a tese
deve permanecer metodológica. Qualquer afirmação de superioridade, robustez
clínica ou generalização externa será adicionada apenas se estiver diretamente
apoiada por resultados, intervalos de incerteza e análise qualitativa.
